\subsection{Error of using SSHC as a proxy for SLC in GMSL estimation}

The key result from the per-source error characterisation shows that the proxy error differs in both sign and magnitude between ice sources: Greenland melt leads to a systematic underestimation of GMSLC, while West and East Antarctic sources produce a smaller bias of opposite sign. This source dependence means that the SSHC proxy error cannot be corrected by a single scalar factor; the correction would itself depend on the unknown spatial distribution of ice change. This matches the findings of \textcite{lickleyBiasEstimatesGlobal2018, al-attarReciprocitySensitivityKernels2023}.

At the global scale the resulting error is modest, on the order of 1-10\% depending on expected value and uncertainty of the GMSLC; however, regional averages exhibit substantially larger and spatially variable biases, since SSHC and SLC differ by the solid-earth deformation and geoid perturbation terms, which are functions of the ice load which is spatially heterogeneous. This motivates a more principled inference framework in which the full fingerprint physics is incorporated into the estimation procedure.


\subsection{Inversion for ice thickness change from ocean altimetry}

The spatial recovery of the simple inversion is good with GMSLC estimates within one standard deviation
%TC:ignore
(e.g. \autoref{fig:simple_inversion_example}: true GMSLC = -5.47 mm; posterior expectation = -5.50 mm, posterior standard deviation = 0.0456 mm),
%TC:endignore
showing that the Bayesian framework is able to effectively use the ocean altimetry data to constrain the ice thickness change. The sensitivity tests suggest that even with poor calibration of the priors, the ocean altimetry is informative enough to constrain the ice thickness change, which is encouraging for the application of this method to real data where the priors are likely to be imperfect. The repeatability of this across multiple inversion suggests that this is not a feature of a single true data source, but rather a general feature of the method.

The improved $z$-scores compared to the proxy method confirm that the Bayesian inversion is an improvement for estimating GMSLC. The proxy's large $z$-scores reflect systematic uncertainty underestimation: the SSHC proxy cannot represent the full SLC response, whereas the Bayesian framework can, producing well-calibrated posteriors. Well-calibrated uncertainty is practically significant: overly confident GMSLC estimates propagate through to sea level budget closure and coastal hazard projections, where they can mask discrepancies and mislead risk assessments. This was confirmed by the number of true GMSLC values falling within one posterior standard deviation: 60.5\% under the Bayesian method compared to 14.6\% under the proxy method (n=261).

The computational cost of the inference is sensitive to the choice of prior covariance amplitude: conjugate gradient iteration counts ranged from 47 (amplitude multiplier 0.5$\times$) to 610 (multiplier 5.0$\times$), while length-scale and mean-offset variations had comparatively little effect on convergence.

\subsection{Joint inversion using ice, firn, and ocean dynamic changes}

The joint inversion show good recovery of the detail and good sensitivity when using all data types. There are clear trade-offs between the model parameters, as expected, and the method does well to disentangle these variable sources. The lower recovery of firn in some regions is expected due to low sensitivity of mass sensitive data sources to firn changes and could be masked by the more dominant ice changes.

Knock-out tests show that there is information gained from each data type, especially ice altimetry constraining the trade-off between the different densities of firn and ice effectively. SSH altimetry is critical for constraining ocean dynamic topography, which makes intuitive sense as it provides the spatial coverage needed to separate ocean dynamic signals from ice-driven fingerprints.

The loss in ocean dynamic topography detail is acceptable, given the altimetry noise parameter (1\,mm standard deviation) and so only recovering features over this noise level is expected. This is also in part a consequence of jointly inverting for multiple sources, which introduces additional degrees of freedom which lead to some detail loss.
